\subsection{Seeded defects: what testing misses and the proofs reject}
\label{sec:mutation}
\begin{table}[tp]
\centering
\caption{Seeded defects in the extracted optimizer and exporter. T1: number of the 380 runs (296 on the corpus of 148 generated automata under two export policies, 84 on the benchmark patterns) whose output differs from the unmutated tool. T2: index (from zero) of the first random word on which the language test finds a difference; a dash means none was found. Fuzz: inputs on which the language test finds a mismatch, over inputs on which the mutant's output differs from the unmutated tool's, among 3,000 random automata (5,000 words per output). Outcome: ``caught'' by T2; ``missed by T1, T2'' but found by the fuzz; ``missed by all tests'' (M12, for which a hand-built automaton shows the language change); ``no change found'' (M05 is a different composition of verified passes, so its language is preserved by the pass theorems; M08, M21 and M22 are argued equivalent). Rocq: all 22 transcriptions are rejected (Table~\ref{tab:rocq-rejection}).}
\label{tab:mutants}
\scriptsize
\setlength{\tabcolsep}{3pt}
\begin{adjustbox}{max width=\textwidth}\begin{tabular}{@{}lp{0.34\textwidth}rrrp{0.13\textwidth}l@{}}
\toprule
Mutant & Seeded defect (pass) & T1 & T2 first word & Fuzz & Outcome of the tests & Rocq \\
\midrule
M01 & closed bound 0 on a just-reset clock & 0 & \textemdash{} & 44/234 & missed by T1, T2 & rejected \\
M02 & invariant from min, not max, of branch bounds & 30 & 44 & 154/428 & caught & rejected \\
M03 & x$\ge$2, x$\le$2 always empty (strictness ignored) & 120 & 0 & 1940/3456 & caught & rejected \\
M04 & lower-bound tests ignored for clock liveness & 104 & 0 & 918/1650 & caught & rejected \\
M05 & normalization before dead-reset removal omitted & 78 & \textemdash{} & 0/1116 & no change found & rejected \\
M06 & guard $x>0$ treated as trivially true & 38 & 0 & 770/1284 & caught & rejected \\
M07 & clocks merged without equal reset sets & 14 & 0 & 10/32 & caught & rejected \\
M08 & entry into initial location not excluded & 0 & \textemdash{} & 0/4 & no change found & rejected \\
M09 & accepting and non-accepting locations merged & 6 & 0 & 56/116 & caught & rejected \\
M10 & resets ignored in simulation test & 0 & \textemdash{} & 2/4 & missed by T1, T2 & rejected \\
M11 & parallel transitions merged without comparing resets & 62 & 1 & 184/850 & caught & rejected \\
M12 & complement of $x<c$ taken as $x>c$ & 0 & \textemdash{} & 0/0 & \textbf{missed by all tests} & rejected \\
M13 & entry lower bound from max, not min & 88 & 0 & 1616/2098 & caught & rejected \\
M14 & split copies entered the wrong way round & 15 & 0 & 93/195 & caught & rejected \\
M15 & entry condition not updated for resets & 40 & 0 & 263/462 & caught & rejected \\
M16 & reset clock with $x\le0$ judged false & 0 & \textemdash{} & 31/82 & missed by T1, T2 & rejected \\
M17 & implication direction in disjunct pruning & 1 & \textemdash{} & 5/29 & missed by T1, T2 & rejected \\
M18 & invariants ignored when merging exported locations & 8 & 84 & 36/328 & caught & rejected \\
M19 & initial location pruned as unreachable & 376 & 0 & 6000/6000 & caught & rejected \\
M20 & $x\le u$ taken to imply $x<w$ for $u\le w$ & 24 & 40 & 43/220 & caught & rejected \\
M21 & safety check of dead-reset pass removed & 0 & \textemdash{} & 0/0 & no change found & rejected \\
M22 & weaker bound for conjunctive guard & 0 & \textemdash{} & 0/84 & no change found & rejected \\
\bottomrule
\end{tabular}\end{adjustbox}

\end{table}
\begin{table}[tp]
\centering
\caption{The same mistakes applied to the Rocq definitions: lemma on which the development first fails to compile, with the time in seconds from the start of the build of the mutated file to the first error (Rocq 9.0.1). The first five rows are the defects that T1 and T2 missed (M12 also escapes the fuzz); the others follow in numerical order. M05, M08, M21 and M22 do not change the language, and are rejected because the proofs name the mutated step or need the removed condition.}
\label{tab:rocq-rejection}
\small
\begin{tabular}{@{}llr@{}}
\toprule
Mutant & First failing lemma & Time (s) \\
\midrule
M01 & \texttt{inv\_guard\_w\_eq} & 4.7 \\
M10 & \texttt{sim\_spec} & 3.8 \\
M12 & \texttt{compl\_cases} & 4.8 \\
M16 & \texttt{const\_true\_spec} & 3.0 \\
M17 & \texttt{insert\_g\_holds} & 4.1 \\
\midrule
M02 & \texttt{max\_all\_ge} & 9.1 \\
M03 & \texttt{contradictory\_bounds\_sound} & 2.7 \\
M04 & \texttt{rename\_item\_holds} & 3.2 \\
M05 & \texttt{optimize\_step\_accepts} & 5.4 \\
M06 & \texttt{trivial\_c\_sound} & 3.8 \\
M07 & \texttt{mergeable\_b\_sound} & 2.8 \\
M08 & \texttt{mergeable\_b\_sound} & 3.3 \\
M09 & \texttt{states\_mergeable\_fast\_eq} & 4.3 \\
M11 & \texttt{same\_ends\_spec} & 4.6 \\
M13 & \texttt{min\_all\_le} & 1.8 \\
M14 & \texttt{notlonger\_complete} & 2.5 \\
M15 & \texttt{split\_trans\_in} & 3.1 \\
M18 & \texttt{dmerge\_states\_ext} & 4.7 \\
M19 & \texttt{dprune\_accepts} & 7.7 \\
M20 & \texttt{implies\_c\_sound} & 2.2 \\
M21 & \texttt{remove\_dead\_resets\_clock\_fast\_eq} & 3.5 \\
M22 & \texttt{ub\_guard\_sound} & 1.9 \\
\bottomrule
\end{tabular}

\end{table}
The agreement above shows that the tools are correct on the sample, not what the proofs add. To measure that, we seeded 22 plausible defects, one at a time, in a scratch copy of the extracted optimizer and exporter (Table~\ref{tab:mutants}; for instance a strict instead of non-strict test at the bound zero, a reset set not compared when two locations are merged, the complement of $x<c$ taken as $x>c$). Each mutant was tested in three ways: T1 compares its output with that of the unmutated tool on the benchmark and the corpus of 148 generated automata; T2 tests, for up to 100{,}000 random words per input, whether the output automaton accepts the words of its input; and the fuzz repeats this on 3{,}000 random automata drawn with the same generator and another seed (a corpus 25 times larger, with small constants including zero), testing 5{,}000 words on each output that differs from the unmutated tool's. Only outputs that differ are tested, because an output byte-identical to the unmutated tool's has the language of the unmutated tool, which the proofs equate with the input's.

Eighteen of the 22 mutants change the language of some automaton; four do not (M05 changes the output but not the language; M08, M21 and M22 change neither). Thirteen changes are found by T2. Five, M01, M10, M12, M16 and M17, are found by neither T1 nor T2: the benchmark and the small corpus never reach the defect (four of the five change no output at all), and the word test, which fails to notice even a single dropped transition in 36\% of our controls, does not hit the few words that separate the languages. The fuzz finds M01, M10, M16 and M17, on 44, 2, 31 and 5 automata of 3{,}000; the one mismatch for M17 is a single word, so the earlier classification of M17 as language-preserving, which rested on T2 alone, was wrong. M12 changes no output on any of the 3{,}000 automata. For M01 and M12 we built an automaton on which the mutant's language differs (for M01, 852 accepted sample words become none; for M12, 3{,}234 against 3{,}564), so M12 is a real language change that every test of this kind missed. The test oracles are themselves unverified code, and neither an exact language-equivalence check nor targeted word generation was tried.

We then applied the same mistakes to the Rocq definitions in a scratch copy of the development (control: the unmutated copy compiles in about 43 s with the four standard axioms). All 22 are rejected by Rocq within 1.8 to 9.1 seconds at a named lemma (Table~\ref{tab:rocq-rejection}). The rejection is at the lemma that ties the implementation to its specification (for M01, against the specification function of the invariants), so the check is only as strong as that specification: a mutation that changed implementation and specification together would pass. It is also conservative: it rejects the four mutants that do not change the language, so a rejection means that the definition no longer matches the specification, not that its language differs. The Rocq mutants are transcriptions of the extracted-code mutants, made by the authors. These results make a bounded claim: for these defects in this code, T1 and T2 missed five of 18 language-changing defects, a fuzz 25 times larger missed one, and the development rejects all of them. They do not show that CASAAL contains such defects, and CASAAL may well be correct.


